\chapter{Машина без диригента}
% full version: chapters/02_glutcomputer.tex

У вашому комп'ютері є диригент --- тактовий генератор. Мільярди разів на секунду він дає відмашку, і всі елементи схеми перемикаються одночасно, в ногу. У мозку жодного диригента немає. Кожен нейрон працює сам по собі: сигнали приходять до нього коли доведеться і йдуть від нього коли доведеться. Погляньмо, що можна обчислити в такій машині, якщо взяти лише \nt{GLUT}-нейрони --- самі «плюси».

\section*{Діряве відро}

Нейрон зручно уявити відром із діркою. Кожен імпульс, що надійшов, --- це ківш води. Вода потроху витікає крізь дірку. Якщо ковші надходять часто, вода встигає піднятися до краю --- і відро перекидається: нейрон клацає, відро спорожніло, усе спочатку. Якщо ковші надходять рідко, вода встигає витекти, і нічого не відбувається.

\pencilfig{2}{\pencilart{2}}{Нейрон --- діряве відро: краплі входів наповнюють його, вода витікає; дійшла до риски --- клацання.}

Дірка --- головна відмінність нейрона від логічного вентиля. Нейрон пам'ятає імпульс, що надійшов, не вічно, а лише кілька мілісекунд. Усе, що прийшло в межах цього вікна, додається; те, що було раніше, забуто.

\section*{«Або» та «і»}

Якщо одного ковша досить, щоб перекинути відро, нейрон спрацьовує на будь-який вхід. Це «або». Якщо потрібні два ковші, починається цікаве: «чи прийшли обидва?» без диригента не має сенсу, доки не сказано --- \emph{у межах якого часу}. Нейрон спрацює, лише якщо два імпульси прийшли майже одночасно, поки перший ківш не витік. Це «і в часі», детектор збігів. В одних нейронів вікно --- кілька мілісекунд, у нейронів слуху --- частки мілісекунди.

\section*{Чого самі «плюси» не вміють}

Спробуйте побудувати «ні» --- нейрон, який працює, коли вхід мовчить. Не вийде: тиша не наливає води у відро. І це правильно не лише для одного нейрона, а й для будь-якої мережі з самих «плюсів»: якщо входи посилити, жоден нейрон у такій мережі не почне працювати менше. Звідси найнеприємніша властивість: \emph{активність не можна погасити активністю}. Пожежу, що розгорілася, можна зупинити лише одним способом --- перестати підкидати дрова.

Природа знайшла обхідний шлях. У багатьох органах чуттів той самий подразник передають дві групи нейронів. В оці одні клітини повідомляють «стало світліше», інші --- «стало темніше». Якщо замість одного проводу у вас пара проводів «так» і «ні», то «ні» виходить задарма: просто поміняйте проводи місцями. З такими парами з самих «плюсів» можна зібрати будь-яку логічну схему --- і навіть зрозуміти, що обчислення завершено: у кожній парі засвітився рівно один провід. Так інженери будують асинхронні схеми, які правильно працюють за будь-яких затримок.

\section*{Час із затримок}

Без диригента немає спільного годинника, але час можна обчислювати й так. Звук від джерела збоку приходить у ближнє вухо раніше, ніж у дальнє, --- на частки мілісекунди. Прокладімо від лівого вуха провід зліва направо вздовж ряду детекторів збігів, а від правого --- справа наліво. Два сигнали побіжать назустріч один одному й зустрінуться в тій точці, куди обидва добіжать одночасно. Спрацює детектор саме в цьому місці. Різниця в часі перетворилася на \emph{номер} нейрона, що спрацював, --- тобто на напрямок на звук. Точність виходить у десяток мікросекунд, хоча кожен окремий нейрон значно менш точний: рятує те, що детекторів багато.

\pencilfig{102}{%
% delay line: the left-ear wire runs right along the top, the right-ear wire runs left along the bottom
\draw[pen=pcBlue] (0,0.6) -- (5.6,0.6); \draw[pen=pcBlue] (6.0,-0.6) -- (0.4,-0.6);
\foreach \i in {1,...,7}{
  \pgfmathsetmacro{\x}{0.75*\i}
  \draw[pen=pcBlue] (\x,0.6) -- (\x,0.24); \draw[pen=pcBlue] (\x,-0.6) -- (\x,-0.24);
  \ifnum\i=5 \draw[dotpen=pcAmber, wash=pcAmber, line width=1.1pt] (\x,0) circle (0.2);
  \else \draw[dotpen=pcBlue, wash=pcBlue] (\x,0) circle (0.2); \fi}
\draw[arr=pcBlue] (0.95,0.78) -- (1.75,0.78); \draw[arr=pcBlue] (5.3,-0.78) -- (4.5,-0.78);
\node[lbl, anchor=east] at (-0.05,0.6) {ліве вухо};
\node[lbl, anchor=west] at (6.05,-0.6) {праве вухо};
\node[lbl, text=pcAmber!70!black, anchor=south] at (3.75,0.95) {тут зустрілися};
\draw[arr=pcAmber!70!black] (3.75,0.95) -- (3.75,0.68);
% sound source on the left
\draw[penlite] (-1.25,-0.25) -- (-1.05,-0.25) -- (-0.8,-0.45) -- (-0.8,0.05) -- (-1.05,-0.15) -- (-1.25,-0.15) -- cycle;
\foreach \r in {0.2,0.35}{\draw[penlite] ($(-0.75,-0.2)+(-40:\r)$) arc (-40:40:\r);}
\node[lbl, anchor=north] at (-0.95,-0.5) {звук зліва};
}{Звук зліва раніше приходить у ліве вухо, його сигнал устигає пробігти далі, і зустріч виходить праворуч від середини. Номер нейрона, що спрацював, --- напрямок на звук.}

\section*{Пам'ять --- це вогнище}

Візьмімо групу нейронів, які сильно збуджують одне одного. Якщо її розпалити коротким сигналом, вона підтримуватиме сама себе, як вогнище, що саме підкидає собі дрова. Сигнал скінчився, а група все горить --- вона пам'ятає, що сигнал був. Це комірка пам'яті. А якщо сигнал був закороткий, вогнище не встигає розгорітися й гасне.

Але ось біда: як стерти таку пам'ять? Активністю --- ніяк. Лишається чекати, доки нейрони втомляться і вогнище прогорить.

\section*{Таймер, запущений подією}

З'єднаймо групи в ланцюжок: перша запалює другу, друга третю. Штовхнімо першу --- і ланцюжком побіжить хвиля, як кісточками доміно. Номер кісточки, що впала, показує, скільки минуло часу від поштовху. Це не годинник, що цокає завжди, а таймер, який вмикається подією і відпрацьовує один раз. У співочих птахів під час пісні нейрони одного відділу спрацьовують суворо по черзі --- дуже схоже на такий ланцюжок.

Отже, з самих «плюсів» виходить чимало. Бракує трьох речей: вибрати одне з багатьох, стерти пам'ять за командою і вберегти мережу від пожежі. Усі три дають «мінуси» --- про них наступний розділ.

\fullversion{2}
